\section{Phát biểu vấn đề}
\label{sec:problem}

\subsection{Khoảng trống nghiên cứu}

Hầu hết bài so sánh LTSF, kể cả bảng kết quả của LTSF-Linear \cite{zeng2023ltsf},
chỉ nói model nào thắng trên từng bộ dữ liệu và từng horizon.
Cách đó giúp xếp hạng, nhưng \textbf{chưa giúp người dùng chọn model} khi gặp chuỗi mới:

\begin{quote}
      Với một chuỗi chưa từng gặp, nên dùng phân rã (DLinear) hay chuẩn hóa bằng giá trị cuối
      (NLinear)? Dựa vào đâu để chọn?
\end{quote}

Toner và Darlow còn chỉ ra rằng trên benchmark ``sạch'', DLinear và Linear thường gần như
giống nhau \cite{toner2024linear}. Khi sai số trung bình gần nhau như vậy, chỉ nhìn bảng MSE
không đủ: cần xem \textbf{đặc trưng của từng chuỗi} (xu hướng mạnh hay yếu, mùa rõ không,
có không dừng hay không) mới giải thích được vì sao DLinear hoặc NLinear thắng.

Các chỉ số độ mạnh xu hướng và mùa \cite{hyndman2021fpp3,cleveland1990stl} và hướng xử lý
không dừng / dịch phân phối như RevIN \cite{kim2022revin} đã có sẵn trong tài liệu,
nhưng ít bài gắn chúng vào so sánh DLinear với NLinear rồi rút ra
\textbf{một quy tắc chọn model ngắn, dùng lại được}.
Đó chính là khoảng trống nhóm muốn lấp.


\subsection{Câu hỏi nghiên cứu}

\begin{description}[style=nextline,leftmargin=2.2em]
      \item[RQ1: Đặc trưng dữ liệu đến mô hình thắng.]
            Những đặc trưng đo được nào (độ mạnh xu hướng \(F_T\), độ mạnh mùa \(F_S\),
            và proxy không dừng / dịch mức) dự đoán DLinear hay NLinear đạt lỗi thấp hơn
            dưới cùng một protocol?
      \item[RQ2: Lookback và đặc trưng dữ liệu.]
            Độ dài lookback \(L\) làm thay đổi thứ hạng tương đối của NLinear và DLinear như thế nào,
            và pattern độ nhạy đó có phụ thuộc profile đặc trưng dữ liệu không?
\end{description}

\textbf{Đóng góp mục tiêu (các tuần sau).}
Một scatter hoặc heatmap giữa đặc trưng dữ liệu và khoảng cách hiệu năng
\(\Delta = \mathrm{MSE}(\mathrm{DLinear}) - \mathrm{MSE}(\mathrm{NLinear})\),
cộng một đoạn guideline ngắn cho người thực hành.
Giả thuyết sai vẫn là kết quả hợp lệ nếu protocol được cố định và công bố rõ.

\subsection{Bộ dữ liệu}

Để có phương sai về xu hướng, mùa và tính dừng, nhóm dùng các bộ LTSF chuẩn cộng thêm
một chuỗi giá cổ phiếu (miền tài chính):

\begin{center}
      \begin{tabular}{@{}ll@{}}
            \toprule
            \textbf{Dataset} & \textbf{Vai trò}                                                      \\
            \midrule
            ETTh1, ETTh2     & Máy biến áp điện (theo giờ)                                           \\
            Weather          & Biến khí tượng                                                        \\
            Exchange-Rate    & Tỷ giá; thường không dừng mạnh                                        \\
            Electricity      & Tải điện nhiều chiều                                                  \\
            \textbf{VIC}     & Giá đóng cửa ngày cổ phiếu Vingroup (HOSE), \(\log(\mathrm{close})\);
            dùng để kiểm tra lại trên dữ liệu tài chính                                              \\
            \bottomrule
      \end{tabular}
\end{center}

VIC \emph{không} phải benchmark duy nhất: nó bổ sung một chuỗi nhiễu, dễ đổi trạng thái thị trường,
ngoài các CSV LTSF thông thường, để guideline không chỉ đúng trên kiểu điện / thời tiết.
